%% Bounded editorial treatment of planned 77.1 and the elementary-model core of 77.2.
\chapter{Proper Forcing and Master Conditions}
\label{sec:p06-proper-forcing-master}

A master condition makes the part of each dense set visible in a countable
elementary submodel usable below every further strengthening. It is not just
one choice compatible with one member of each dense set.

The comparison survey in \cref{sec:p06-forcing-variants} introduces the
preceding forcing vocabulary and presentations. Here the elementary-model
criterion is developed for individual posets, with positive constructions and
a failed master-condition test. The cited forcing infrastructure of
Chapter~\ref{ch:constructibility-forcing} remains a separate input, not a
consequence of that criterion. Iteration and forcing-axiom constructions need
additional arguments beyond the single-poset scope of
\cref{def:p06-pfm-boundary}.

Work in classical \(\mathsf{ZFC}\) as metatheory. Fix a nonempty forcing
poset \(\mathbf P=(P,\leq_P,1_P)\), where \(1_P\) is greatest and
\(r\leq_P q\) means stronger. Compatibility, ccc, density and poset filters
retain \cref{def:p06-ccc,def:p06-poset-filter}. The countable \(N\) below
is usually \emph{not transitive}; it is not the externally countable
transitive ground model of the earlier pedagogical forcing construction.

The elementary-model criterion follows selected passages of
\cite[pp.~19--23,48--49]{Dow1988ElementarySubmodels}; the distinct
ordinal-name criterion is
\cite[Definition~6.6, pp.~44--45]{VialeAudritoSteila2014SemiproperV1}.
The local proofs expand the constructions, without claiming either entire
source or Shelah's original book.

\section{The ambient set and countable elementary models}
\label{sec:p06-pfm-models}

\begin{hbdefinition}{def:p06-pfm-hchi}{Hereditary-size ambient set}
For a set \(x\), \(\operatorname{tc}(\{x\})\) is the least transitive set
containing \(x\) as a member. For regular uncountable \(\chi\), put
\[
 H_\chi=\{x:|\operatorname{tc}(\{x\})|<\chi\}.
\]
For the forcing tests fix
\[
 \kappa=\max(\aleph_1,|\operatorname{tc}(\{\mathbf P\})|),
 \qquad \chi\text{ regular},\quad\chi>2^\kappa.
\]
This is an \emph{admissible height for this treatment}, a size convention,
not recursion-theoretic admissibility or inaccessibility.
Write \(N\prec(H_\chi,\in)\) for an elementary membership substructure
with countable underlying set \(N\), and require \(\mathbf P\in N\).
The components \(P,\leq_P,1_P\) then belong to \(N\), by their unique
definitions as projections of that finite tuple.
\end{hbdefinition}

\begin{hblemma}{thm:p06-pfm-hchi}{proved=annex:p06.pfm.hchi}{The hereditary-size closure needed here}
The class \(H_\chi\) is a transitive set. Every subset of any \(a\in H_\chi\)
belongs to \(H_\chi\); finite tuples of its members belong to it.
For nonempty countable \(a\in H_\chi\), a surjection \(\omega\to a\)
belongs to \(H_\chi\). At the declared heights, \(\omega_1,\mathbf P\in
H_\chi\). No Power Set axiom for \(H_\chi\) is used.
\end{hblemma}

\begin{hbdefinition}{def:p06-pfm-hulls}{Witness functions and hulls}
Fix a well-order \(\triangleleft\) of \(H_\chi\). For each existential
membership formula \(\exists y\,\varphi(y,\bar x)\), let \(F_\varphi\)
return the least witness in \(H_\chi\), when one exists, and a fixed default
otherwise. These are countably many finite-arity functions.
For a countable seed \(A\subseteq H_\chi\) containing the parameters, put
\[
 A_0=A,\qquad A_{n+1}=A_n\cup
 \{F_\varphi(\bar a):\bar a\text{ is a matching-arity tuple from }A_n\},
 \qquad \operatorname{Hull}(A)=\bigcup_{n<\omega}A_n .
\]
Further specified countably many finite-arity operations and constants
may be included.
\end{hbdefinition}

\begin{hbtheorem}{thm:p06-pfm-hulls}{proved=annex:p06.pfm.hulls}{Hulls, countable elements and ordinal traces}
The hull is countable and elementary in \((H_\chi,\in)\).
The countable sets closed under the specified operations form a club on
\([H_\chi]^\omega\), the countably infinite subsets: unbounded under inclusion and closed under unions
of increasing countable chains. Any two clubs in this sense have club
intersection.
For countable \(N\prec(H_\chi,\in)\), \(\omega\subseteq N\), and every countable
\(a\in N\) is a subset of \(N\). When \(\chi>\omega_1\),
\(\delta_N=N\cap\omega_1\) is a countable limit ordinal below \(\omega_1\).
An arbitrary uncountable element of \(N\) need not be a subset of \(N\).
\end{hbtheorem}

The actual witness condition of \cref{thm:tarski-vaught}, not closure under
a few familiar set operations, establishes elementarity. The distinction
between \(a\in N\) and \(a\subseteq N\) is essential in the ccc argument.

\section{Predensity and the master-condition test}
\label{sec:p06-pfm-master}

\begin{hbdefinition}{def:p06-pfm-predensity}{Predensity below a condition}
For \(X\subseteq P\), define
\[
 X^\downarrow=\{s\in P:\exists x\in X\ (s\leq_P x)\}.
\]
It is \emph{predense below \(q\)} iff
\[
 \forall r\leq_P q\ \exists x\in X\ \exists s\in P\
                        (s\leq_P r\land s\leq_P x).
\]
Write \(r\perp_P x\) for incompatibility. Predensity requests compatibility,
not \(x\leq_P r\), and not membership of \(q\) in \(X\).
\end{hbdefinition}

\begin{hbdefinition}{def:p06-pfm-master-condition}{Generic or master condition for a model}
A condition \(q\in P\) is \emph{\((N,\mathbf P)\)-generic}, or a
\emph{master condition for \(N\)}, iff \(D\cap N\) is predense below \(q\)
for every dense \(D\subseteq P\) with \(D\in N\).
There is no requirement \(q\in N\).
Density and compatibility are metatheoretic; their bounded poset assertions
agree with those in the transitive \(H_\chi\) containing the complete datum.
\end{hbdefinition}

\begin{hblemma}{thm:p06-pfm-predensity}{proved=annex:p06.pfm.predensity}{Dense closure and filter interpretation}
Predensity of \(X\) below \(q\) is equivalent to density of
\(X^\downarrow\) below \(q\). The master test is unchanged if only dense
\emph{open} sets in \(N\) are tested.
If \(V\) is a transitive ground model of \(\mathsf{ZFC}\), the data belong
to \(V\), and \(G\) is a \(V\)-generic poset filter in an outer universe
containing \(q\), then predensity implies \(G\cap X\ne\varnothing\).
If predensity fails, some \(r\leq_P q\) is incompatible with every \(x\in X\);
no filter containing \(r\) can meet \(X\).
Generic-filter existence is not inferred from a vacuous quantification.
\end{hblemma}

\begin{hblemma}{thm:p06-pfm-antichains}{proved=annex:p06.pfm.antichains}{Maximal-antichain form}
Every dense \(D\subseteq P\) contains a maximal antichain of \(P\).
If \(D,\mathbf P\in N\), one such antichain belongs to \(N\).
The master test is equivalent to requiring predensity of \(A\cap N\)
below \(q\) for every maximal antichain \(A\in N\).
\end{hblemma}

\begin{hbdefinition}{def:p06-pfm-proper}{Properness and its height quantifiers}
The poset is \emph{proper} iff, for every admissible regular \(\chi\),
every countable \(N\prec(H_\chi,\in)\) containing \(\mathbf P\), and every
\(p\in P\cap N\), there is a master condition \(q\leq_P p\).
This tests all declared sufficiently large heights, not just one selected
model or height. The stationary-set-preservation characterization is not
used as a premise.
\end{hbdefinition}

\section{Two complete properness constructions}
\label{sec:p06-pfm-constructions}

\begin{hbtheorem}{thm:p06-pfm-ccc}{proved=annex:p06.pfm.ccc}{Countable and ccc forcings are proper}
If \(P\) is countable, or more generally \(\mathbf P\) is ccc, every \(q\in P\)
is a master condition for every model in the properness test.
In particular \(q=p\) suffices.
\end{hbtheorem}

\begin{hbdefinition}{def:p06-pfm-countable-closure}{Countable closure}
Every descending sequence \(p_{n+1}\leq_P p_n\), \(n<\omega\), has a common
lower bound in \(P\). This includes sequences in the metatheory, not just
sequences which belong to \(N\).
\end{hbdefinition}

\begin{hbtheorem}{thm:p06-pfm-closed}{proved=annex:p06.pfm.closed}{Countably closed forcings are proper}
Assume \(\mathbf P\) is countably closed. Every \(p\in P\cap N\) has a master strengthening: construct a descending
sequence inside \(N\) meeting its dense sets, then take a common lower
bound in \(P\). That bound need not belong to \(N\), and the construction
also works for dense sets which are not open.
\end{hbtheorem}

\begin{hbtheorem}{thm:p06-pfm-examples}{proved=annex:p06.pfm.examples}{The sufficient hypotheses are different}
Cohen forcing \(2^{<\omega}\), stronger toward end-extension, is countable
and proper but not countably closed. Its generic union is a new binary real.
Countable partial functions \(\omega_1\to2\), stronger toward extension,
give a countably closed proper poset which is not ccc.
Properness is therefore not a no-new-reals assertion or a second name for
either sufficient hypothesis.
\end{hbtheorem}

\section{The ordinal bound in semigenericity}
\label{sec:p06-pfm-semiproper}

Bounds in this section are ground-model ordinals. The recursive forcing
theorem, including ordinal decision and truth/density, is imported forcing
infrastructure, not a newly proved theorem here. Its cited boundary in the
earlier forcing chapter remains. In particular, this is not an assumption
that the forcing relation is absolute between \(H_\chi\) and the universe.
When a ground model \(V\) is displayed, it is transitive and satisfies
\(\mathsf{ZFC}\); \(\omega_1,\chi,N\) in that forcing discussion are
computed in \(V\), and generic extensions are taken in an outer universe.
The preceding combinatorial theorems are ZFC results and do not require
the existence of an externally countable transitive model.

\begin{hbdefinition}{def:p06-pfm-semigeneric}{Semigenericity and the club test}
A condition \(q\) is \emph{\((N,\mathbf P)\)-semigeneric} iff, for every
\(P\)-name \(\dot\alpha\in N\cap H_\chi\) with
\(1_P\forces\dot\alpha<\check\omega_1\), one has
\[
 q\forces\dot\alpha<\check{\delta_N}.
\]
The \emph{semiproperness test at \(\chi\)} requires a club
\(C\subseteq[H_\chi]^\omega\) of elementary models containing \(\mathbf P\)
such that every \(p\in P\cap N\), \(N\in C\), has a semigeneric
strengthening. The uniform-height convention requires the test at every
admissible regular \(\chi\); its club may depend on the height.
This is a club-model ordinal-name test, not the universal-model demand
on every dense set in the definition of properness.
\end{hbdefinition}

\begin{hbtheorem}{thm:p06-pfm-semigeneric}{proved=annex:p06.pfm.semigeneric}{Properness supplies the semigeneric club test}
Using the imported ordinal-decision/truth interface, every proper poset
satisfies the semiproperness test at every admissible height.
At each height expand the membership structure by operations which record
a deciding maximal antichain and its ordinal-value map for each relevant
ordinal name. For countable elementary models of that expansion, every
master condition is semigeneric.
\end{hbtheorem}

No converse, one-height/global equivalence, iteration theorem or general
stationary-preservation characterization follows from this implication.

\section{A failed master-condition test}
\label{sec:p06-pfm-failure}

\begin{hbtheorem}{thm:p06-pfm-collapse}{proved=annex:p06.pfm.collapse}{Finite collapse has no master strengthening}
Let \(P\) consist of finite partial functions \(\omega\to\omega_1\),
ordered by reverse inclusion, with empty function greatest.
For every model \(N\) in the properness test and \(q\in P\), choose
\(n\notin\operatorname{dom}(q)\) and set \(r=q\cup\{(n,\delta_N)\}\).
The dense \(D_n=\{s\in P:n\in\operatorname{dom}(s)\}\) belongs to \(N\),
but \(r\) is incompatible with all of \(D_n\cap N\).
The poset is not proper and also fails the semiproperness club test:
on a club containing its explicit ordinal-name coding ledger, no condition
is semigeneric. A generic union is a surjection from \(\omega\) onto the
ground \(\omega_1\).
\end{hbtheorem}

\begin{hbwarning}{def:p06-pfm-weak-test}{One compatible choice is insufficient}
In every poset, elementarity gives \(d\in D\cap N\) below \(p\) whenever
\(p\in P\cap N\) and \(D\in N\) is dense.
A weakened test requesting only one member of \(D\cap N\) compatible with
\(p\) would therefore pass every poset, including the finite collapse.
The missing quantifier is \(\forall r\leq_P q\).
\end{hbwarning}

\begin{proofboundary}{def:p06-pfm-boundary}{Single-poset scope}
These complete elementary-model and single-poset constructions do not
supply semiproper iteration, PFA or MM consistency, or the full
Shelah 1982 and Abraham 2010 source presentations. Countable-support
preservation additionally uses the later lifted-model, prescribed-initial
and parameter-invariance constructions in
Sections~\ref{sec:p06-proper-iteration-lemma}
and~\ref{sec:p06-properness-parameter-invariance}.
\end{proofboundary}
